\clearpage
\section{Especificação formal M12.2: roteiros, compartilhamento e Storage}
\label{sec:formal-m12-2}
\sloppy

M12.2 formaliza o alvo normativo de \textit{Roteiro\_Experimento},
compartilhamento entre professores, upload/download via Storage e associação
opcional a Post (RF21--RF23, UI-11, Q09, Q13), na composição com M7
(idempotência), M9 (autorização no commit), M11 (vínculo canônico) e a fronteira
abstrata \texttt{roteiro\_anexo}/\texttt{acessoRoteiroValidado} de M12.1. Ela
refina aditivamente a fronteira de M12.1 sem reabrir sua prova executável. A
evidência cobre somente o modelo formal declarado.

\subsection*{CUE}
O contrato CUE (\texttt{\#M12\_2Contrato}) verifica shapes de Roteiro
(provisório/validado/publicável), referência canônica ao objeto no Storage,
snapshot imutável do anexo e sua ligação explícita à referência canônica,
compartilhamento único por \texttt{(roteiro, professor)}, contexto de
autorização/usuário, vínculo canônico, contexto de Post, projeção mínima de
download, URL já emitida com validade temporal abstrata e operação idempotente.
Os limites são nomes 1--150, geração 1--120, tamanho positivo e PDF
\textbf{estritamente menor} que 15 MiB (15\,728\,640 bytes). O snapshot do anexo
embute \texttt{\#M12\_1RoteiroAnexo} e acrescenta \texttt{geracao}: o contrato
M12.1 antigo \textbf{não} prova a geração; a prova é dada por M12.2. CUE não
prova autorização nem concorrência.

\subsection*{Dono, upload e publicabilidade}
O dono \texttt{id\_professor\_upload} é imutável e coincide com \texttt{owner\_uid}.
O cadastro cria um Roteiro \texttt{PROVISORIO} sem geração; a validação do objeto
fixa a geração persistida a partir da geração real do objeto e move para
\texttt{VALIDADO}; publicar move para \texttt{PUBLICAVEL}. O objeto e a geração
são fixos (\texttt{ObjetoNuncaSobrescrito}, \texttt{GeracaoNaoMudaSemValidar}) e
o Roteiro só é elegível a compartilhamento/anexo/download quando publicável.
Firestore e Storage não têm transação única: a geração é persistida no documento
e comparada na transação do Firestore; a inspeção do objeto no Storage é etapa
externa; a janela residual é coberta pela imutabilidade do objeto vinculado e
pela reconciliação de órfãos.

\subsection*{Compartilhamento, revogação Q09 e download}
Só o proprietário compartilha ou revoga, sempre para professor ativo e distinto,
com no máximo uma relação por \texttt{(roteiro, professor)}
(\texttt{CompartilhamentoUnico}). A revogação impede novas associações
(\texttt{RevogadoNaoEmite}), mas \textbf{preserva} os Posts históricos e seus
snapshots. O professor acessa por propriedade ou compartilhamento atual; o aluno,
pelo vínculo canônico atual e um Post acessível que referencia o roteiro; o
Chefe, por Q13; um \textbf{ex-aluno} é negado ainda que a claim esteja atualizada
(\texttt{WitnessExAlunoNaoEmite}). Turma \texttt{Arquivada} é somente leitura:
membro com vínculo atual ainda baixa o anexo de Post histórico, mas nenhuma
escrita acadêmica é aceita (\texttt{TurmaArquivadaNegaEscritaRoteiro});
Post removido não autoriza (\texttt{PostRemovidoNegaAluno}).

A \textbf{URL já emitida} é modelada como fronteira temporal abstrata
(\texttt{urlsAtivas}): a remoção, a perda de vínculo e a revogação impedem
\textbf{novas emissões} (\texttt{EmissaoExigeAcesso}), mas uma URL ativa
permanece utilizável até expirar (\texttt{UrlAtivaUsavel},
\texttt{WitnessUrlAtivaAposRevogacao}), sem prometer revogação retroativa. A
geração da URL vinculada é a geração canônica.

\subsection*{Anexo a Post, histórico e composição M7/M9}
No máximo um anexo por Post. Anexar, trocar e \emph{manter} o anexo exigem
acesso atual ao roteiro no commit, \textbf{inclusive na edição}; desvincular não
exige acesso. O snapshot guarda a geração validada
(\texttt{AnexoEhGeracaoValidada}, \texttt{AnexoUsaGeracaoCanonica}); trocar e
desvincular registram o snapshot anterior no histórico, que nunca é removido
(\texttt{HistoricoAnexoNuncaRemovido}); remover o Post da apresentação preserva
anexos e histórico (\texttt{RemocaoPostPreservaSnapshot}). A composição M7
modela primeira execução (grava receipt sob o \texttt{idOperacao}), retry
idêntico (não reaplica) e reuso incompatível (rejeitado); M9 faz a revogação de
vínculo impedir o commit posterior.

\subsection*{Responsabilidade das camadas e limites}
CUE valida shapes e limites. Alloy valida posse, transições, ACL de UID único,
revogação Q09, turma arquivada, ex-aluno, Post removido, anexo/histórico e
composição M7/M9. Rust valida proveniência, ordem exata do receipt e as regras
determinísticas de limite estrito de 15 MiB, tipo, titularidade, status,
ligação do anexo à referência canônica, emissão/uso de URL, aluno/ex-aluno e
unicidade de UID. A evidência \textbf{não} certifica \texttt{functions/},
frontend, Firestore/Storage Rules, Auth, Storage real, bytes binários,
atomicidade entre serviços nem concorrência.

\begingroup\small\raggedright
% Gerado por lcqui-spec-doc; não editar.
\subsection*{M12\_2\_Roteiros\_Compartilhamento\_Storage --- formalização executável M12.2}
Shapes estruturais de Roteiro (provisório/validado/publicável), referência canônica ao objeto no Storage com geração, snapshot imutável do anexo refinando aditivamente \#M12\_1RoteiroAnexo, ligação explícita do anexo à referência canônica, compartilhamento único por (roteiro, professor), contexto de autorização (M9), vínculo canônico (M11), contexto de Post, projeção mínima de download, URL já emitida com validade temporal abstrata, operação idempotente M7 e efeito mínimo de notificação. CUE verifica campos, enums fechados, nulabilidade, condicionais e limites (nome 1..150, tamanho 1..15728639, geração 1..120, PDF strictly < 15 MiB). Dono imutável, upload provisório->validado->publicável, objeto/geração fixos, compartilhamento único, revogação Q09, ACL de professor versus aluno via vínculo atual + Post acessível, Chefe Q13, turma arquivada somente leitura, ex-aluno negado mesmo com claim atualizada, Post removido, anexar/trocar/manter/desvincular, histórico imutável, distinção entre nova emissão de URL e uso da URL já emitida (sem revogação retroativa) e composição M7/M9 são verificados em Alloy; limites e canonicalização determinística por Rust. CUE não prova autorização, concorrência nem atomicidade Firestore/Storage.
\begin{description}
\item[id\_roteiro] Tipo: string. Obrigatório: sim. Aceita nulo: não.
SQL: TEXT.
Identidade do Roteiro; dono imutável e referência canônica.
\item[id\_professor\_upload] Tipo: string. Obrigatório: sim. Aceita nulo: não.
SQL: TEXT.
Dono imutável; coincide com owner\_uid do objeto.
\item[nome] Tipo: string. Obrigatório: sim. Aceita nulo: não.
SQL: TEXT.
Nome de apresentação do Roteiro; de 1 a 150 caracteres.
\item[storage\_path] Tipo: string. Obrigatório: sim. Aceita nulo: não.
SQL: TEXT.
Referência canônica ao objeto no Storage; nunca credencial.
\item[content\_type] Tipo: string. Obrigatório: sim. Aceita nulo: não.
SQL: TEXT.
Formato validado; PDF (application/pdf) na V1.
\item[tamanho\_bytes] Tipo: int. Obrigatório: sim. Aceita nulo: não.
SQL: INTEGER.
Tamanho validado; PDF estritamente menor que 15 MiB (15728640 bytes).
\item[owner\_uid] Tipo: string. Obrigatório: sim. Aceita nulo: não.
SQL: TEXT.
Titular do objeto; coincide com id\_professor\_upload.
\item[geracao] Tipo: string. Obrigatório: sim. Aceita nulo: não.
SQL: TEXT.
Geração/versão validada e imutável do objeto; persistida no Firestore.
\item[status\_roteiro] Tipo: enum. Obrigatório: sim. Aceita nulo: não.
SQL: ENUM.
Valores: PROVISORIO, VALIDADO, PUBLICAVEL.
Upload provisório não é publicável; publicável após objeto validado e geração fixada.
\item[nome\_arquivo] Tipo: string. Obrigatório: sim. Aceita nulo: não.
SQL: TEXT.
Nome do arquivo no snapshot imutável do anexo; até 150 caracteres.
\item[id\_operacao] Tipo: string. Obrigatório: sim. Aceita nulo: não.
SQL: TEXT.
Identidade de comando M7 (uid, tipo\_operacao, payload\_hash).
\item[tipo\_operacao] Tipo: enum. Obrigatório: sim. Aceita nulo: não.
SQL: ENUM.
Valores: CADASTRAR\_ROTEIRO, VALIDAR\_OBJETO, PUBLICAR\_ROTEIRO, COMPARTILHAR\_ROTEIRO, REVOGAR\_COMPARTILHAMENTO, ANEXAR\_ROTEIRO\_POST, MANTER\_ROTEIRO\_POST, DESVINCULAR\_ROTEIRO\_POST, EMITIR\_URL.
Operação idempotente M7 sobre roteiro, compartilhamento ou anexo a Post.
\item[via] Tipo: enum. Obrigatório: sim. Aceita nulo: não.
SQL: ENUM.
Valores: PROPRIETARIO, COMPARTILHADO, ALUNO\_POST, CHEFE\_Q13.
Caminho de autorização do download; Chefe por Q13 apenas em moderação/auditoria.
\item[validade\_url] Tipo: enum. Obrigatório: sim. Aceita nulo: não.
SQL: ENUM.
Valores: ATIVA, EXPIRADA.
Validade temporal abstrata da URL já emitida; não há revogação retroativa.
\item[revogado\_em] Tipo: string. Obrigatório: sim. Aceita nulo: sim.
SQL: TIMESTAMP.
Instante da revogação do compartilhamento; nulo enquanto vigente.
\item[removido\_da\_apresentacao] Tipo: bool. Obrigatório: sim. Aceita nulo: não.
SQL: BOOLEAN.
Post removido da apresentação; snapshot e histórico preservados (Q09/RF25).
\end{description}

% Gerado por lcqui-spec-doc; não editar.
\subsection*{Evidência formal M12.2 --- roteiros, compartilhamento, Storage/download e anexo a Post}
CUE verifica shapes de Roteiro (provisório/validado/publicável), referência canônica ao objeto no Storage com geração, refinamento aditivo do snapshot do anexo sobre \#M12\_1RoteiroAnexo, ligação explícita do anexo à referência canônica, compartilhamento único por (roteiro, professor), contexto de autorização/usuário, vínculo canônico, contexto de Post, projeção mínima de download, URL já emitida com validade temporal abstrata e operação idempotente, com enums fechados, nulabilidade, condicionais e limites (nome 1..150, tamanho 1..15728639, geração 1..120, PDF estritamente menor que 15 MiB). Alloy verifica dono imutável, upload provisório->validado->publicável, objeto e geração fixos, compartilhamento único e destinatário ativo, revogação Q09 preservando Posts/snapshots históricos, ACL de professor versus acesso do aluno via vínculo canônico atual e Post acessível, Chefe em Q13, turma arquivada somente leitura, ex-aluno negado mesmo com claim atualizada, Post removido, anexar/trocar/manter/desvincular, histórico de anexo imutável, distinção entre nova emissão de URL (revalida acesso) e uso da URL já emitida (prazo abstrato, sem revogação retroativa) e composição M7 (primeira execução grava receipt; retry não reexecuta; reuso incompatível rejeitado) e M9 (revogação de vínculo impede commit). Rust valida proveniência, ordem exata do receipt e as regras determinísticas de limite estrito de 15 MiB, tipo, titularidade, status publicável, ligação do anexo à referência canônica, emissão/uso de URL, aluno/ex-aluno e unicidade de UID. A evidência não certifica Firebase, Storage real, bytes binários, Rules, Auth, atomicidade entre Firestore e Storage, a janela residual de reconciliação nem concorrência sob carga.
\begin{description}
\item[M12\_2-INV-001] CadastroComecaProvisorio: UNSAT.\newline Escopo: \texttt{check CadastroComecaProvisorio for 6}.
\item[M12\_2-INV-002] PublicavelSoDeValidado: UNSAT.\newline Escopo: \texttt{check PublicavelSoDeValidado for 6}.
\item[M12\_2-INV-003] ValidarFixaGeracao: UNSAT.\newline Escopo: \texttt{check ValidarFixaGeracao for 6}.
\item[M12\_2-INV-004] RoteiroValidadoTemGeracao: UNSAT.\newline Escopo: \texttt{check RoteiroValidadoTemGeracao for 6}.
\item[M12\_2-INV-005] DonoImutavel: UNSAT.\newline Escopo: \texttt{check DonoImutavel for 6}.
\item[M12\_2-INV-006] ObjetoNuncaSobrescrito: UNSAT.\newline Escopo: \texttt{check ObjetoNuncaSobrescrito for 6}.
\item[M12\_2-INV-007] GeracaoNaoMudaSemValidar: UNSAT.\newline Escopo: \texttt{check GeracaoNaoMudaSemValidar for 6}.
\item[M12\_2-INV-008] SemObjetoNaoEmite: UNSAT.\newline Escopo: \texttt{check SemObjetoNaoEmite for 6}.
\item[M12\_2-INV-009] AnexoUsaGeracaoCanonica: UNSAT.\newline Escopo: \texttt{check AnexoUsaGeracaoCanonica for 6}.
\item[M12\_2-INV-010] AnexoUnicoPorPost: UNSAT.\newline Escopo: \texttt{check AnexoUnicoPorPost for 6}.
\item[M12\_2-INV-011] CompartilharSoProprietario: UNSAT.\newline Escopo: \texttt{check CompartilharSoProprietario for 6}.
\item[M12\_2-INV-012] CompartilharExigeDestinatarioAtivo: UNSAT.\newline Escopo: \texttt{check CompartilharExigeDestinatarioAtivo for 6}.
\item[M12\_2-INV-013] CompartilhamentoUnico: UNSAT.\newline Escopo: \texttt{check CompartilhamentoUnico for 6}.
\item[M12\_2-INV-014] RevogarSoProprietario: UNSAT.\newline Escopo: \texttt{check RevogarSoProprietario for 6}.
\item[M12\_2-INV-015] NaoProprietarioNaoRevoga: UNSAT.\newline Escopo: \texttt{check NaoProprietarioNaoRevoga for 6}.
\item[M12\_2-INV-016] EmissaoExigeAcesso: UNSAT.\newline Escopo: \texttt{check EmissaoExigeAcesso for 6}.
\item[M12\_2-INV-017] ExAlunoNaoEmite: UNSAT.\newline Escopo: \texttt{check ExAlunoNaoEmite for 6}.
\item[M12\_2-INV-018] InativoNaoEmite: UNSAT.\newline Escopo: \texttt{check InativoNaoEmite for 6}.
\item[M12\_2-INV-019] AlunoDependeDeVinculo: UNSAT.\newline Escopo: \texttt{check AlunoDependeDeVinculo for 6}.
\item[M12\_2-INV-020] PostRemovidoNegaAluno: UNSAT.\newline Escopo: \texttt{check PostRemovidoNegaAluno for 6}.
\item[M12\_2-INV-021] TurmaArquivadaNegaEscritaRoteiro: UNSAT.\newline Escopo: \texttt{check TurmaArquivadaNegaEscritaRoteiro for 6}.
\item[M12\_2-INV-022] UrlEmitidaSobreviveARevogacao: UNSAT.\newline Escopo: \texttt{check UrlEmitidaSobreviveARevogacao for 6}.
\item[M12\_2-INV-023] UrlExpiradaNaoUsavel: UNSAT.\newline Escopo: \texttt{check UrlExpiradaNaoUsavel for 6}.
\item[M12\_2-INV-024] UrlAtivaUsavel: UNSAT.\newline Escopo: \texttt{check UrlAtivaUsavel for 6}.
\item[M12\_2-INV-025] AnexarSoDonoDaTurma: UNSAT.\newline Escopo: \texttt{check AnexarSoDonoDaTurma for 6}.
\item[M12\_2-INV-026] AnexarExigeAcessoAtual: UNSAT.\newline Escopo: \texttt{check AnexarExigeAcessoAtual for 6}.
\item[M12\_2-INV-027] ManterExigeAcessoAtual: UNSAT.\newline Escopo: \texttt{check ManterExigeAcessoAtual for 6}.
\item[M12\_2-INV-028] TrocarExigeAcessoAtual: UNSAT.\newline Escopo: \texttt{check TrocarExigeAcessoAtual for 6}.
\item[M12\_2-INV-029] RemocaoPostPreservaSnapshot: UNSAT.\newline Escopo: \texttt{check RemocaoPostPreservaSnapshot for 6}.
\item[M12\_2-INV-030] HistoricoAnexoNuncaRemovido: UNSAT.\newline Escopo: \texttt{check HistoricoAnexoNuncaRemovido for 6}.
\item[M12\_2-INV-031] HistoricoAnexoPreservaObjeto: UNSAT.\newline Escopo: \texttt{check HistoricoAnexoPreservaObjeto for 6}.
\item[M12\_2-INV-032] AnexoEhGeracaoValidada: UNSAT.\newline Escopo: \texttt{check AnexoEhGeracaoValidada for 6}.
\item[M12\_2-INV-033] ChefeNaoAnexa: UNSAT.\newline Escopo: \texttt{check ChefeNaoAnexa for 6}.
\item[M12\_2-INV-034] ChefeNaoCompartilha: UNSAT.\newline Escopo: \texttt{check ChefeNaoCompartilha for 6}.
\item[M12\_2-INV-035] RevogacaoVinculoImpedeCommit: UNSAT.\newline Escopo: \texttt{check RevogacaoVinculoImpedeCommit for 6}.
\item[M12\_2-INV-036] PrimeiraExecucaoProduzReceipt: UNSAT.\newline Escopo: \texttt{check PrimeiraExecucaoProduzReceipt for 6}.
\item[M12\_2-INV-037] IdentidadeComandoUnica: UNSAT.\newline Escopo: \texttt{check IdentidadeComandoUnica for 6}.
\item[M12\_2-INV-038] ReusoIncompativelRejeitado: UNSAT.\newline Escopo: \texttt{check ReusoIncompativelRejeitado for 6}.
\item[M12\_2-INV-039] RetryNaoReexecuta: UNSAT.\newline Escopo: \texttt{check RetryNaoReexecuta for 6}.
\item[M12\_2-INV-040] RetryNaoDuplicaFato: UNSAT.\newline Escopo: \texttt{check RetryNaoDuplicaFato for 6}.
\item[M12\_2-INV-041] TransicoesPreservamCoerencia: UNSAT.\newline Escopo: \texttt{check TransicoesPreservamCoerencia for 6}.
\item[M12\_2-WIT-042] WitnessTurmaAtiva: SAT.\newline Escopo: \texttt{run WitnessTurmaAtiva for 4}.
\item[M12\_2-WIT-043] WitnessCadastraProvisorio: SAT.\newline Escopo: \texttt{run WitnessCadastraProvisorio for 6}.
\item[M12\_2-WIT-044] WitnessValidaObjeto: SAT.\newline Escopo: \texttt{run WitnessValidaObjeto for 6}.
\item[M12\_2-WIT-045] WitnessPublica: SAT.\newline Escopo: \texttt{run WitnessPublica for 6}.
\item[M12\_2-WIT-046] WitnessProprietarioEmite: SAT.\newline Escopo: \texttt{run WitnessProprietarioEmite for 6}.
\item[M12\_2-WIT-047] WitnessCompartilhadoEmite: SAT.\newline Escopo: \texttt{run WitnessCompartilhadoEmite for 6}.
\item[M12\_2-WIT-048] WitnessAlunoEmitePostAtivo: SAT.\newline Escopo: \texttt{run WitnessAlunoEmitePostAtivo for 6}.
\item[M12\_2-WIT-049] WitnessAlunoEmiteTurmaArquivada: SAT.\newline Escopo: \texttt{run WitnessAlunoEmiteTurmaArquivada for 6}.
\item[M12\_2-WIT-050] WitnessExAlunoNaoEmite: SAT.\newline Escopo: \texttt{run WitnessExAlunoNaoEmite for 6}.
\item[M12\_2-WIT-051] WitnessRevogadoNaoEmite: SAT.\newline Escopo: \texttt{run WitnessRevogadoNaoEmite for 6}.
\item[M12\_2-WIT-052] WitnessUrlAtivaAposRevogacao: SAT.\newline Escopo: \texttt{run WitnessUrlAtivaAposRevogacao for 6}.
\item[M12\_2-WIT-053] WitnessChefeEmite: SAT.\newline Escopo: \texttt{run WitnessChefeEmite for 6}.
\item[M12\_2-WIT-054] WitnessPostRemovidoNaoEmiteAluno: SAT.\newline Escopo: \texttt{run WitnessPostRemovidoNaoEmiteAluno for 6}.
\item[M12\_2-WIT-055] WitnessAnexaComAcesso: SAT.\newline Escopo: \texttt{run WitnessAnexaComAcesso for 6}.
\item[M12\_2-WIT-056] WitnessMantemComAcesso: SAT.\newline Escopo: \texttt{run WitnessMantemComAcesso for 6}.
\item[M12\_2-WIT-057] WitnessDesvinculaSemAcesso: SAT.\newline Escopo: \texttt{run WitnessDesvinculaSemAcesso for 6}.
\item[M12\_2-WIT-058] WitnessTrocaPreservaHistorico: SAT.\newline Escopo: \texttt{run WitnessTrocaPreservaHistorico for 6}.
\item[M12\_2-WIT-059] WitnessRetryAposExecucao: SAT.\newline Escopo: \texttt{run WitnessRetryAposExecucao for 6}.
\item[M12\_2-WIT-060] WitnessReusoIncompativel: SAT.\newline Escopo: \texttt{run WitnessReusoIncompativel for 6}.
\item[M12\_2-WIT-061] WitnessObjetoAusenteNaoEmite: SAT.\newline Escopo: \texttt{run WitnessObjetoAusenteNaoEmite for 6}.
\end{description}
UNSAT para check indica ausência de contraexemplo no escopo declarado; SAT para run indica testemunha encontrada. A evidência não certifica a implementação Firebase/backend.

\par\endgroup

\subsection*{Limites}
Esta evidência prova somente o modelo formal no escopo declarado
(\textit{bounded model checking}) e não é homologação da implementação real.
\fussy
